\chapter{Kinematika a Newtonova dynamika}

Pro urceni pohyboveho stavu hmotniho bodu potrebujeme znat, v nejakem case, polohu a rychlost a vyslednici sil, ktera na hmotny bod pusobi. Pak muzeme spocist vsechen budouci vyvoj polohy hmotneho bodu.

\section{Kinematika hmotného bodu}
Poloha částice v trojrozměrném prostoru s kartezkymi souradenicmi v čase $t$ je popsaná polohovým vektorem $\vect{r}(t) \in \mathbb{R}^3$. V ortonormální kartézské bázi $\{\vect{e}_x, \vect{e}_y, \vect{e}_z\}$ máme vyjádření:
\begin{equation}
    \vect{r}(t) = x(t)\vect{e}_x + y(t)\vect{e}_y + z(t)\vect{e}_z,
\end{equation}
nebo ekvivalentni vyjadreni v inych souradnicovych systemech, cylindricke ($\rho(t)$, $\phi(t)$ a $z(t)$) nebo sfericke ($r(t)$, $\theta(t)$, $\phi(t)$).

\emph{Prumerna} velkost rychlost mezi body 1 a 2, mezi ktorymi urazi hmotni bod drahu $\Delta s$ za cas $\Delta t$ je
$$ \bar{v} = \frac{\Delta s}{\Delta t}. $$

Okamžitá rychlost $\vect{v}(t)$ a zrychlení $\vect{a}(t)$ jsou definovány pomocí časových derivací polohového vektoru:
\begin{equation}
    \vect{v}(t) = \odv{\vect{r}}{t} = \dot{\vect{r}}(t), \qquad
    \vect{a}(t) = \odv{\vect{v}}{t} = \dodv{\vect{r}}{t} = \ddot{\vect{r}}(t),
\end{equation}
t.j., pro male casove zmeny $\delta t$, poloha bodu v case $\vect{r}(t + \delta t) \approx \vect{r}(t) + \vect{v}(t)\delta t$ a podobne pro rychlost $\vect{v}(t)$. V limite $\Delta t, \Delta s \to 0$ mame $|\vect{v}| = \bar{v}$.

Pohyb nazyvame \emph{primocary}, kdyz se deje v primce a \emph{krivocary} pro vsetky ostatni pripady. Podlve velikosti rychlosti mame pohyb \emph{rovnomerny} pro $|\vect{v}|=$~konst. a \emph{nerovnomerny} pro $|\vect{v}|\neq$~konst.

Primocary pohyb s konstantnim zrychlenim $a$
$$ x = \frac{1}{2}a t^2 + v_0 t + x_0, $$
kde $v_0$, $x_0$ jsou pocatecni podminky.

\begin{priklad}[Rovnomerny pohyb po kruznici]
    Nech sa hmotny pod pohybuje po kruznici s polomerem $R$ a periodou (jeden obeh) $T$. Najprirozenejsi zapis je pomoci polarnich souradnic,
    \begin{align*}
        r(t) &= R\\
        \phi(t) &= \frac{2\pi}{T}t = \omega t,
    \end{align*}
    kde sme definovali uhlovou frekvenci $\omega = 2\pi/T$.

    Ekvivalentne muzeme pohyb zapsat v kartezkych souradnicich
    $$ x = R \cos(\omega t),\;\; y = R \sin(\omega t) $$.

    Pro rychlost dostavame
    $$ \vect{v} = \odv{\vect{r}}{t} = R\omega (-\sin(\omega t), \cos(\omega t)), $$
    kde hned vidime, ze rychlost $|\vect{v}| = R\omega$ se nemeni v case a je kolma na $\vect{r}$

    Pro  zrychleni dostavame
    $$ \vect{a} = \odv{\vect{v}}{t} = -R\omega^2 (\cos(\omega t), \sin(\omega t)) = - \omega^2 \vect{r}, $$
    kde vidime, ze zrychleni je \emph{dostredive}, t.j., miry do stredu rotace a je kolme na rychlost.

    Jak by vypadala rychlost a zrychleni pro pohyb po elipse? Je velikost rychlosti stale konstantni, pri konstantni uhlove rychlosti?
\end{priklad}

Pro pohyb po spirale

% \begin{definice}[Hybnost a síla]
% Pro hmotný bod o konstantní hmotnosti $m$ pohybující se rychlostí $\vect{v}$ definujeme \textbf{hybnost}:
% \begin{equation}
%     \vect{p} \coloneqq m \vect{v}.
% \end{equation}
% Působí-li na bod celková síla $\vect{F}$, je časová změna hybnosti rovna této síle.
% \end{definice}

% \section{Newtonovy pohybové zákony}
% Základem newtonovské formulace jsou tři pohybové postuláty.

% \begin{veta}[Druhý Newtonův pohybový zákon]
% Časová změna hybnosti hmotného bodu je přímo úměrná výsledné působící síle $\vect{F}$ a má její směr:
% \begin{equation}
%     \odv{\vect{p}}{t} = \vect{F}(\vect{r}, \dot{\vect{r}}, t).
% \end{equation}
% V případě konstantní hmotnosti $m = \text{konst.}$ přechází rovnice do známého tvaru:
% \begin{equation}
%     m \ddot{\vect{r}} = \vect{F}.
% \end{equation}
% \end{veta}

% \begin{priklad}[Jednorozměrný lineární harmonický oscilátor]
% Uvažujme hmotný bod o hmotnosti $m$ na pružině s tuhostí $k > 0$ bez tlumení. Pohybová rovnice má tvar:
% \begin{equation}
%     m \ddot{x} + k x = 0 \iff \ddot{x} + \omega_0^2 x = 0, \quad \text{kde } \omega_0 = \sqrt{\frac{k}{m}}.
% \end{equation}
% Obecné řešení této lineární diferenciální rovnice druhého řádu je:
% \begin{equation}
%     x(t) = A \cos(\omega_0 t + \varphi_0),
% \end{equation}
% kde amplituda $A \ge 0$ a fázový posuv $\varphi_0 \in [0, 2\pi)$ plynou z počátečních podmínek $x(0)$ a $\dot{x}(0)$.
% \end{priklad}

% \begin{poznamka}
% Při transformaci do neinerciální vztažné soustavy (např. rotující úhlovou rychlostí $\vect{\omega}$) vystupují na pravé straně pohybové rovnice setrvačné síly: odstředivá síla $-m\vect{\omega}\crossprod(\vect{\omega}\crossprod\vect{r})$, Coriolisova síla $-2m(\vect{\omega}\crossprod\vect{v})$ a Eulerova síla $-m(\dot{\vect{\omega}}\crossprod\vect{r})$.
% \end{poznamka}
